\worksheettitle
  {Домашняя практика: ответы}
  {\modulemark{M01-H} Точка, прямая, отрезок и луч}

\begin{teacherbox}[Критерии]
  Проверяются число концов, направление продолжения, порядок букв в луче и
  геометрически корректное построение. Небрежность линии отделяется от ошибки
  в модели объекта.
\end{teacherbox}

\section{Ответы}

\begin{enumerate}
  \item Прямая: концов нет, продолжается в обе стороны. Отрезок: два конца,
    за них не продолжается. Луч: одно начало, продолжается в одну сторону.
  \item a) луч $MN$; б) $CD$, $DC$; в) прямая $PQ$ или $QP$.
  \item Прямая проходит через $A,B$ и продолжается за обе точки; отрезок
    заканчивается в $A,C$; луч начинается в $C$, проходит через $B$.
  \item Причина: первая буква ошибочно принята за проходящую точку. Луч
    $PK$ начинается в $P$ и проходит через $K$.
\end{enumerate}

\begin{resultbox}[Маршрут]
  Если ошибки повторяются в заданиях 1, 2 или 4, назначьте M01-R. Если модель
  верна, но построение неточно, достаточно дополнительного построения с
  линейкой.
\end{resultbox}
